\section{The Quantum Simulator}
\label{sec:quantum-simulator}

\subsection{Statevector representation}

The simulator follows the standard statevector formulation~\cite{nielsenchuang2010}. An
$n$-qubit register (QuantumCraft supports $n \in \{1,2,3\}$) is represented by a complex
amplitude vector $\psi \in \mathbb{C}^{2^n}$ indexed by basis state, using the little-endian
convention that qubit~0 is the least-significant bit of the index -- so, for two qubits, the
index ordering is $|00\rangle, |01\rangle, |10\rangle, |11\rangle$ with qubit~0 varying
fastest. A freshly constructed register is initialized to $\psi = (1, 0, \ldots, 0)$, the
all-zero basis state $|0\cdots0\rangle$, and every gate application below preserves the
invariant $\sum_i |\psi_i|^2 = 1$ by construction, since every matrix applied is unitary.

\subsection{Single-qubit gates}

A single-qubit gate is specified by a $2\times2$ unitary matrix; the mod exposes six presets,
each cycled by right-clicking a gate block (Section~\ref{sec:architecture}):
\begin{equation}
X = \begin{bmatrix} 0&1\\1&0 \end{bmatrix}, \quad
Y = \begin{bmatrix} 0&-i\\i&0 \end{bmatrix}, \quad
Z = \begin{bmatrix} 1&0\\0&-1 \end{bmatrix}, \quad
H = \frac{1}{\sqrt2}\begin{bmatrix} 1&1\\1&-1 \end{bmatrix}, \quad
S = \begin{bmatrix} 1&0\\0&i \end{bmatrix}, \quad
T = \begin{bmatrix} 1&0\\0&e^{i\pi/4} \end{bmatrix}.
\end{equation}
Applying such a matrix to qubit $q$ of an $n$-qubit register means: for every pair of basis
indices $(i, j)$ differing only in bit $q$ (with bit $q$ clear in $i$ and set in $j$), replace
$(\psi_i, \psi_j)$ with the matrix product of the $2\times2$ gate and the two-component vector
$(\psi_i, \psi_j)$, leaving every other amplitude untouched -- the direct implementation of
the tensor-product action $I \otimes \cdots \otimes U \otimes \cdots \otimes I$ without ever
materializing the full $2^n \times 2^n$ matrix.

\subsection{Controlled gates and the Toffoli}

A controlled-$U$ gate, of which CNOT ($U = X$) is the special case this mod's control/target
marker pair represents (Section~\ref{sec:architecture}), applies the single-qubit matrix $U$
to a target qubit only for those basis states in which a nominated control qubit's bit is~1,
leaving every basis state with the control bit clear entirely unchanged. This is implemented
as one general mechanism parameterized by an arbitrary $2\times2$ matrix, so a controlled-$X$
(CNOT), controlled-$Z$, or controlled-$Y$ all reduce to the same code path rather than three
separate implementations. The Toffoli gate (CCX), the doubly-controlled generalization
represented by two control markers and one target marker sharing a time column
(Section~\ref{sec:architecture}), applies the $X$ matrix's basis-flip only to those states in
which \emph{both} nominated control qubits' bits are~1: for every pair of basis indices
differing only in the target bit, with both control bits set in each, the two amplitudes are
swapped. Because a classical bit can be encoded as the basis state $|0\rangle$ or
$|1\rangle$ and never placed in superposition, the Toffoli gate applied to classical inputs
reproduces exactly the reversible-AND truth table Toffoli's own reversible-computing
construction describes~\cite{toffoli1980reversible}, verified directly in
Section~\ref{sec:validation} and demonstrated as a worked experiment in
Section~\ref{sec:experiments}.

\subsection{Measurement}

A measurement (triggered by a shift-right-click of the probe, Section~\ref{sec:architecture})
implements the standard Born rule~\cite{nielsenchuang2010}: a uniform random number in
$[0,1)$ is drawn and compared against the cumulative sum of $|\psi_i|^2$ over basis states in
index order, and the first index at which the cumulative sum exceeds the draw is the sampled
outcome. Every amplitude is then set to zero except the sampled outcome's, which is set to
$1$, collapsing the register exactly as an ideal projective measurement would. Because this
sampling procedure is only ever invoked once per probe use and the resulting register is
immediately discarded rather than persisted (Section~\ref{sec:architecture}), no mechanism is
needed to reset a register's state between uses, unlike a physical quantum computer or a
notebook-based simulator session, where explicit re-preparation between shots is often a
separate step.

\subsection{On-demand circuit execution}

A \texttt{QuantumCircuit} is simply an ordered list of operations -- single-qubit gates,
controlled gates, and Toffoli gates -- built directly from the world scan
(Section~\ref{sec:architecture}) rather than persisted as block-entity state anywhere. Running
a circuit allocates a fresh register initialized to $|0\cdots0\rangle$ and applies every
recorded operation to it in the scanned left-to-right (time) order, then discards the
intermediate register once a final result -- an ideal readout or a single measurement -- has
been reported to the player. There is, deliberately, no notion of a circuit's state persisting
between one probe use and the next: every probe use is a fresh, independent run of the whole
circuit from $|0\cdots0\rangle$, matching how a real quantum circuit is executed once as a
complete unit rather than incrementally advanced tick by tick the way CircuitCraft's
continuous MNA solve is~\cite{picos2026circuitcraft}.
